\section{Conclusions}
\label{iii:sec:conclusiones}
\begingroup
\fontsize{10.5}{12.8}\selectfont
\setlength{\parskip}{2pt plus .5pt minus .5pt}

The central result is a reversible constitutive determination on
the contractive domain: the two oriented channels of a transport
produce four related coordinates, and those coordinates recover
the angular pair. Structure first fixes the response equation;
evaluation determines the values while preserving the orientation
needed to reconstruct their spectral antecedents. This two-way
direction, constructive and inverse, gives precise content to the
thesis of structure, equation and value.

\subsection*{The response preserves the angular pair}

The incidences $90$ and $120$ determine a unique solution of the
completion equation for the fixed transport. Its scalar function is
a strictly decreasing bijection from $(0,1)$ onto $(3/4,1)$.
Consequently, the cubic equation~\eqref{iii:eq:cubic} has exactly
one admissible root for each channel response. Inversion is not a
numerical coincidence: it reconstructs the two contractions and,
through their logarithms, the angular pair.

The four vacuum coordinates satisfy
\[
\widehat\mu\widehat\varepsilon=\widehat c^{-2},\qquad
\widehat\mu/\widehat\varepsilon=\widehat Z^2.
\]
Within the fixed product, quotient and positivity rules, these
identities determine the quadratic evaluations. Their common
dependency also appears under sheet interchange: permittivity and
permeability are permuted, impedance is inverted and speed is
preserved. The symmetric product and oriented quotient thus transmit
complementary information. Recovery preserves the sheet labels and
internal normalization; the source state additionally preserves
route, boundary and memory.

The factorization $\mathcal V=a\exp(\eta\mathcal R)$ identifies
these two kinds of information within the operator itself. The positive
factor $a=\sqrt{\det\mathcal V}$ determines normalized speed;
the unimodular component determines impedance. Preserving both
reconstructs the response. This separation explains the relation between
unit-product reciprocal sheets and the four vacuum readings:
normalization separates the common mode, which remains necessary
to recover the complete operator.

The inverse links responses evaluating different functions of the
same channels. The ratio of their exponents reconstructs the
anisotropy of the action ellipse and the relative LC impedance,
as proved in Proposition~\ref{iii:prop:forma-LC}. The action
section and the frequency and impedance references complete their dimensional quantities.
The identity $c_{\rm int}=c_{\rm vac}$ in turn establishes that
the kinematic and constitutive readouts realize the same speed
section in a common chart. The normalized trace preserves this
readout under nonadic amplification.

\subsection*{Composition and reconstruction across sectors}

The transported law $r\odot t=f(\psi(r)\psi(t))$ expresses transport
composition in constitutive coordinates. Its identity is $3/4$ in the
boundary extension, and $-\log\psi(r)/30$ converts composition into
addition. Common projectors preserve the operator realisation and sheet
labels. The velocity of a single state, in turn, retains the ordinary
product of its two channels. The two operations are distinguished by
their arguments and by the exact proof of their relationship.

Higgs--impedance reconstruction establishes another explicit connection.
The gauge couplings recover the common contraction and the
electromagnetic coupling; the Higgs self-coupling recovers the oriented
coordinate. Within the tree-level realisation, these data determine
the two channels and their vacuum responses. Domain inequalities and
genealogical compatibility specify which triples belong to the
realisation under consideration. Retaining marked projectors also
allows reconstruction of the operator; the enriched history retains
its additional provenance data.

\subsection*{The cycle relates response, moment and orientation}

The dodecaphasic realization explains the provenance of the rational
function through a determinant quotient in dimensions $3$ and $4$.
The rank difference produces $-1/12$ as a normalized trace on this
finite domain. The quarter-turn plane preserves other information:
its operator satisfies $J^2=-I$, and its resolvent coefficient
retains the orientation sign. The two logarithmic integrations of
that coefficient, with their initial conditions, determine the
Catalan moment at unlimited depth.

The trace table $b_n$ reconstructs $-\log f(s)$ and produces convergent
moments for $\operatorname{Re}z>0$, with an explicit remainder bound.
Its value at depth one is $\log(4/3)$, consistent with the constitutive
endpoint. The Mellin representation and the distinction between diagonal
and tensorial products retain the operations preceding the scalar
evaluation of each moment.

The integral reconstruction in
Theorem~\ref{iii:thm:restitucion-funcional-catalan} completes the
functional correspondence in both directions:
\[
 \mathcal C_2(s)=\int_0^s\log\frac{s}{t}\,
       \frac{1+t}{1+t+t^2}f(t)\,dt,\qquad
 f(s)=\frac{1+s+s^2}{s(1+s)}\mathcal D^2\mathcal C_2(s).
\]
The condition $\mathcal C_2(s)\to0$ at the origin removes all
additive and logarithmic freedom. The endpoint value $G_{\rm Cat}$,
the oriented kernel, and the constitutive response are thus placed
in a single collection with explicit reconstruction operations.
The arguments $s_\pm$ retain their angular determination; the
normalized trace preserves the speed of the same state when linking
it to action and length. This correspondence specifies the structural
content preserved in passing from the collection to its values.

The Catalan--Pell identity specifies an algebraic evaluation of this
collection:
\[
\mathcal C_2(2-\sqrt3)
=\frac23G_{\rm Cat}-\frac{\pi}{12}\log(2+\sqrt3).
\]
The coefficient $2/3$ arises from the calendar's odd classes and
the depth weight; the first two logarithmic derivatives at this section take the values
$\pi/12$ and $1/4$. Reciprocity of the Pell unit links expansion
and contraction, while the two constitutive arguments remain
determined by the angular pair. This relation keeps the collection,
its readouts and the sections of evaluation explicit.

\Needspace{9\baselineskip}
The Barbero--Immirzi functional is subsequently expressed as an
oriented difference of the same function of the reconstructed
channels:
\[
\gamma=\mathcal H(\psi(r_-))-\mathcal H(\psi(r_+)).
\]
Incidence and dodecaphasic return determine its terms; constitutive
inversion transports them to the vacuum response. With the
orientation mark fixed, channel interchange reverses the sign of
$\gamma$; joint conjugation preserves the functional. This
distinction proves that its content exceeds an unlabelled spectrum
and is preserved under normalized amplifications.

% BEGIN SINTESIS_ADITIVA_20260922
\par
Section~\ref{delta22:iii:restitucion} assembles the spectral potential of constitutive response, global channel recovery through normalized speed and Barbero--Immirzi, and stability bounds. On the generated image, retaining labels and bases, the composition recovers the angular pair, action section, and periodic reading of the dodecaphasic record.
% END SINTESIS_ADITIVA_20260922

Harmonic resolution of the dodecaphasic pattern preserves exactly
modes $3,9,4,8$, with unnormalized coefficients $72,72,-6,-6$.
The magnitude ratio $12:1$ expresses the oriented multiplicities.
The incidence defect $-1/12$ and singular coefficient $1/24$
remain assigned to their respective operations. Explicit comparison
with the equal-weight reader, whose coefficient is $1/1080$,
specifies how the oriented chain enters the normalization of return
while preserving the complete functionals beyond their poles.

\subsection*{From arithmetic balance to electrical relations}

The vacancy discrepancy provides a centred polarization whose
magnitude remains below one unit and whose increment satisfies an
exact temporal balance. Changes between floor and ceiling
conventions are recorded by a boundary term. A charge section and
an oriented duration realize that balance as centred charge and
current, preserving its telescoping law.

The electronic prefactor likewise preserves its antecedent
structure. Enumeration of the fibres in $D_9^3$ determines the
tritic singular mode of value $1/2$; the complete spectrum of their
pairing proves
\[
\|[P_{\Sigma,3},P_{\Pi,3}]\|=\sqrt3/4.
\]
Its quadratic defect $3/4$ coincides with the limit of the
$90/120$ sector quotient. This equality relates results of
specified operators, each with its own space and action: the
prefactor is an arithmetic incompatibility norm, whereas the
interior response depends on its contractive channels.

The positive charge section determined by
$Z_{\rm vac}e_a^2=2\alpha h_a$ combines coupling, action and
impedance. The von Klitzing, conductance, flux and Josephson
relations satisfy exactly
\[
R_KG_0=2,\qquad G_0Z_{\rm vac}=4\alpha,\qquad
K_{J,a}\Phi_{0,a}=1,\qquad K_{J,a}^2R_K=4/h_a.
\]
Action return preserves resistance and conductance; charge and
flux transform with $R_{\rm act}^{1/2}$ and Josephson with
$R_{\rm act}^{-1/2}$. These dependencies jointly explain which
quantities retain the action and which preserve only its
constitutive quotient.

Dimensional realization preserves all these identities when the
action, charge and speed bases change. The resulting dimensionless
values belong to the construction's internal chart; metrological
recognition requires independent evaluation of the dimensional
sections and the corresponding physical regime. The established
mathematical scope is the positive, invertible response, its
compositions and covariance on the stated domains.

Together, these results show how a single APP--TRIT--TPK genealogy
preserves the relations required to pass from paths and arithmetic
operations to channels, moments, constitutive response and electrical
relations. The values thereby acquire an explained provenance and
recoverable information: they express the operations producing
them and reveal, within their representation, the structure
retained in them.
\par\endgroup
